\section{Discussion}
\subsection{Defending the Operational Performance Shift}
The transition from a high-recall baseline regime (Dice: 0.8531) to a high-precision operational regime (Precision: 99.66\%, Dice: 0.7922) is an intentional design choice aimed at high-reliability maritime surveillance. In operational maritime response, the economic and logistical cost of a False-Positive (FP)—mobilizing containment vessels and uncrewed response teams to a non-hazardous slick—is prohibitively higher than the risks associated with marginally lower pixel-wise recall. By prioritizing the "deployment reliability" of the 99.66\% precision threshold, the PGDE architecture effectively eliminates the operational noise that typically hampers autonomous environmental alert systems.

\subsection{Architectural Rationale for Precision Gains}
The 99.66\% precision is specifically achieved through the integrated auxiliary suppression head. Unlike standard CNNs that rely solely on intensity-based contrast, the PGDE model exploits second-order physical signatures. The auxiliary head identifies a signature as a spill only when the Laplacian derivative field confirms the presence of a viscous boundary damping signature characteristic of mechanical oil. This mechanism allows the model to "gate out" lookalike features that share identical backscatter intensity but lack the physical structural complexity of crude oil. Activation heatmaps confirm that the neural focus is anchored on these physical boundary tokens, validating the physics-guided injection as the primary driver for false-alarm suppression.

\subsection{Phase-Specific Training and System Stability}
The two-phase curriculum is foundational for aligning ImageNet-inspired backbones with the stochastic spectral distributions of SAR data. Standard end-to-end training without initial stabilization leads to gradient misalignment, as the encoders over-fit to speckle-induced noise. By anchoring the GGFM to stabilized ResNet-18 features in Phase 1, we established a latent foundation that allowed for efficient global refinement in Phase 2. This strategy ensures the system remains robust across the variable sea states fundamental to coherent sensing.

The integration of PIECL represents a fundamental transition from a purely data-driven segmenter to a physically-constrained surveillance system. While the V9 Dual-Encoder provide the neural "intelligence" to identify dark patches in low-contrast SAR data, the PIECL provides the "reasoning" necessary to resolve ambiguities through surface-roughness and texture density. The PIECL enforces strict spatial constraints ensuring that predicted spill regions are confined to physically valid water surfaces, eliminating false activations on maritime vessels and coastal infrastructure. By enforcing the domain-constraint that mechanical oil spills lack the high-frequency structural signatures of metallic vessel hulls, the system achieves a level of operational robustness that is unattainable through neural architectures alone. The constraint layer does not alter the learned representation, but enforces physically valid outputs during deployment. This hybrid approach ensures that the reported 99.66\% precision is physically grounded and deployment-ready for autonomous global monitoring. Unlike conventional post-processing heuristics, the PIECL is derived from physically interpretable constraints and remains fully decoupled from model training, ensuring that performance gains are attributable to deployment-stage reasoning rather than learned bias.
